\documentclass[a4paper,11pt]{article}

\usepackage[T1]{polski}
\usepackage[utf8]{inputenc}
\usepackage{xcolor}
\usepackage{listings}
\usepackage{minted}
\usepackage[margin=8mm]{geometry}
\setlength{\parindent}{0pt}

% czcionka do kodów w środku tekstu
% Ustawienia listingu
\lstdefinestyle{custom}{
    backgroundcolor=\color{white},
    inputencoding=utf8,
    extendedchars=\true,
    commentstyle=\color{red},
    keywordstyle=\color{blue},
    breaklines=true,
    basicstyle=\ttfamily\footnotesize,
    tabsize=1,
    escapechar='
}
\lstset{style=custom}

\begin{document}

\begin{itemize}
  \item lstlisting: 
  \begin{lstlisting}
    Przyklad kodu
      Przyklad kodu
  \end{lstlisting}

  \begin{lstlisting}[language=java]
    public class PrzykladKodu {
        // Komentarz
        public static void main(String[] args) {
            System.out.println("Przyklad kodu w java");
            System.out.println("Przyklad kodu w java");
        }
    }
  \end{lstlisting}

  \item minted:
  \begin{minted}{java}
    public class PrzykladKodu {
        // Komentarz
        public static void main(String[] args) {
            System.out.println("Przyklad kodu w java");
            System.out.println("Przyklad kodu w java");
        }
    }
  \end{minted}

  \begin{minted}{bash}
    #!/bin/bash
    echo "Przyklad kodu w bash"
    echo "Przyklad kodu w bash"
  \end{minted}
\end{itemize}

\end{document}